% SPDX-License-Identifier: Apache-2.0
\section{A Foreign Module as a Unit}
\label{sec:x-foreign}

A design is seldom all one language. A module that already exists in
Verilog or VHDL, a vendor's or a colleague's, should sit in a TxHDL
design as a unit and be simulated with it, not rewritten.
\code{verilog\_unit()} in the build does that for Verilog: Verilator
compiles the module into a C++ model, a small tool writes a C shim
over the model and the Rust that makes a \code{Unit} of it, an
\code{In} or \code{Out} per port and an \code{Rx} or \code{Tx} per
\code{x\_data}, \code{x\_valid}, \code{x\_ready} trio, and the
crate named holds that unit. \code{vhdl\_unit()} does it for VHDL
with nvc as the engine, a child process behind a testbench that
takes a step per line on its standard input, since nvc is a program
and not a library. Once per rising edge the unit gives the module
its inputs as the edge left them, lets it settle, makes the
transfers the module agrees to on its channels, pulses its clock and
reads its outputs: a registered unit like any other, which goes
first in the join order like any unit with wires out. The runtime
document~\cite{runtime} states the engines' four calls.

The example co-runs the stage of Section~\ref{sec:x-stage} with its
own Verilog and its own VHDL. The build lowers the stage to both,
Verilator makes a unit of the Verilog and nvc one of the VHDL, and
three pipelines run side by side on the same offers, the Rust stage
in one and the two foreign units in the others; the three sinks
must take the same words at the same ticks, and the run asserts
they did. That is the lowering checked live, against the runtime, in
both languages, rather than by replaying a trace as
Section~\ref{sec:x-lower} does. Fig.~\ref{fig:x-foreign} shows the
three pipelines in step.

\begin{figure}[htbp]
\centering
\input{foreign_timing.tex}
\caption{The stage, its Verilog and its VHDL, side by side: the
same words in on \code{inp}, \code{vinp} and \code{hinp}, the same
words out, plus one, on \code{out}, \code{vout} and \code{hout},
at the same ticks.}
\label{fig:x-foreign}
\end{figure}

\lstinputlisting[caption={\code{ex\_foreign.rs}. Run with \code{bazel run //lib/examples:ex\_foreign}.},label={lst:x-foreign}]{lib/examples/ex_foreign.rs}

\lstinputlisting[style=ascii,caption={What \code{ex\_foreign} printed when the build ran it: each word as the three sides took it.},label={lst:x-foreign-out}]{out_foreign.txt}
